\documentclass[a4paper,11pt]{article}

\usepackage[T1]{fontenc}
\usepackage{lmodern}
\usepackage[margin=2.5cm]{geometry}
\usepackage{hyperref}
\usepackage{xcolor}
\usepackage{listings}
\usepackage{booktabs}
\usepackage{enumitem}
\usepackage{fancyvrb}
\usepackage{metalogo}

\hypersetup{
  colorlinks=true,
  linkcolor=blue!60!black,
  urlcolor=blue!70!black,
  citecolor=green!50!black,
  pdftitle={confprogbook — A LaTeX class for conference program books},
  pdfauthor={John Aoga},
}

\lstset{
  language=[LaTeX]TeX,
  basicstyle=\ttfamily\small,
  keywordstyle=\color{blue!70!black}\bfseries,
  commentstyle=\color{green!50!black}\itshape,
  stringstyle=\color{red!60!black},
  breaklines=true,
  columns=fullflexible,
  keepspaces=true,
  frame=single,
  backgroundcolor=\color{gray!5},
  rulecolor=\color{gray!40},
  xleftmargin=1em,
  xrightmargin=1em,
  aboveskip=1em,
  belowskip=1em,
  moretexcs={conferencetitle,conferencesubtitle,conferenceedition,conferencedate,
    conferencelocation,conferenceshort,booktitle,editors,institution,
    institutionaddress,copyrightnotice,conferencelogo,pdfscale,
    makefrontcover,makefrontmatter,cpbbeginprogram,cpbendprogram,
    cpbbeginabstracts,cpbendabstracts,cpbbeginparticipants,cpbendparticipants,
    cpbbegincomments,cpbendcomments,programpartindex,
    dayheading,pleheading,sesheading,puttalk,puttalktitle,putplen,putpart,
    includepaper,programskip,euros,RED,
    setdaycolors,setplenarycolors,setsessioncolors,settableaccentcolor,
    setlandscapecolwidth,setpartname,accentcol,rr},
}


\newcommand{\pkg}[1]{\textsf{#1}}
\newcommand{\cls}[1]{\textsf{#1}}
\newcommand{\cmd}[1]{\texttt{\textbackslash#1}}
\newcommand{\meta}[1]{\ensuremath{\langle}\textit{#1}\ensuremath{\rangle}}
\newcommand{\marg}[1]{\texttt{\{}\meta{#1}\texttt{\}}}
\newcommand{\oarg}[1]{\texttt{[}\meta{#1}\texttt{]}}

\title{%
  \cls{confprogbook} --- A \LaTeX{} class for\\
  conference program books\\[1ex]
  \large Version 1.0 --- 2026/03/11%
}
\author{John Aoga\\\texttt{john.aoga@uclouvain.be}}
\date{}

\begin{document}

\maketitle

\begin{abstract}
\noindent
The \cls{confprogbook} class provides a complete, ready-to-use solution for typesetting
conference \emph{Book of Abstracts} or \emph{Program Books}. It is built on the standard
\cls{book} class and offers high-level commands for cover pages, programmatic tables of
contents with dotted leaders, session and talk entries, PDF abstract inclusion,
participant lists, and organizational comments --- all with customizable colors.
The class was originally developed for the Benelux Meeting on Systems and Control
and has been generalized for use by any academic conference.
\end{abstract}

\tableofcontents
\clearpage

%% ================================================================
\section{Introduction}
\label{sec:intro}
%% ================================================================

Organizing a conference often requires producing a \emph{Book of Abstracts}: a printed
or PDF booklet that contains the program schedule, all submitted paper abstracts
(as included PDFs), a list of participants, and practical information.

The \cls{confprogbook} class encapsulates years of experience producing such books
into a single, easy-to-use \LaTeX{} class. Instead of manually configuring dozens
of packages and defining custom commands, you simply write:

\begin{lstlisting}
\documentclass[a4paper]{confprogbook}
\conferencetitle{10th International Symposium on Examples}
\conferencedate{June 1--3, 2027}
\conferencelocation{Zurich, Switzerland}
\end{lstlisting}

\noindent and the class takes care of page layout, headers/footers, colors,
and provides all the commands you need.

\subsection{Features}

\begin{itemize}[nosep]
  \item Automatic cover page and front matter generation.
  \item Colored, boxed headings for days, plenary talks, and sessions.
  \item Programmatic table of contents with dotted leaders and page references.
  \item Simple command to include paper PDFs with auto-labeling.
  \item Participant list formatting with email \texttt{mailto:} links.
  \item Landscape program overview tables with accent colors.
  \item Fully customizable color scheme.
  \item Compatible with pdf\LaTeX{}, \LuaLaTeX{}, and \XeLaTeX{}.
\end{itemize}

\subsection{License}

This work is released under the \LaTeX{} Project Public License v1.3c or later.
See \url{https://www.latex-project.org/lppl/}.

%% ================================================================
\section{Quick Start}
\label{sec:quickstart}
%% ================================================================

Below is a minimal working example that produces a complete conference program book.

\begin{lstlisting}
\documentclass[a4paper]{confprogbook}

%% --- Conference metadata ---
\conferencetitle{10th International Symposium on Examples}
\conferencedate{June 1--3, 2027}
\conferencelocation{Zurich, Switzerland}
\conferenceshort{ISE 2027}
\booktitle{Book of Abstracts}
\editors{Alice Smith, Bob Jones, and Carol Lee}
\institution{ETH Zurich}
\institutionaddress{Zurich, Switzerland}
\conferencelogo[0.15]{logo.png}

\begin{document}

%% --- Front matter ---
\makefrontcover
\makefrontmatter

%% --- Program ---
\cpbbeginprogram
%% Example program content for confprogbook
%% Replace with your actual conference program

\dayheading{Monday, June 1, 2027}

\programskip

\pleheading{Plenary}{Main Hall}{Welcome \& Opening}{}{Alice Smith}{09:00--09:15}

\programskip

\pleheading{Plenary}{Main Hall}{Advances in Optimal Control}{John Doe}{Alice Smith}{09:15--10:15}

\programskip

\sesheading{MonA01}{Main Hall}
{Machine Learning for Control}
{Carol Lee}{10:30--12:30}

\programskip

\puttalk{MonA01-1\hfill 10.30-10.50}
{\puttalktitle{Deep Reinforcement Learning for Robotic Manipulation}{paper001}}
{Jane Researcher}{University of Example}
{Bob Collaborator}{Institute of Technology}
{}{}
{Jane Researcher}

\puttalk{MonA01-2\hfill 10.50-11.10}
{\puttalktitle{Neural Network Verification for Safety-Critical Systems}{paper002}}
{Sam Author}{Technical University}
{}{}
{}{}
{Sam Author}

\sesheading{MonA02}{Room B}
{Optimization Methods}
{Bob Jones}{10:30--12:30}

\programskip

\puttalk{MonA02-1\hfill 10.30-10.50}
{\puttalktitle{Distributed Optimization in Multi-Agent Systems}{paper003}}
{Alex Writer}{Research Lab}
{Pat Coauthor}{University of Science}
{}{}
{Alex Writer}

\programskip

\dayheading{Tuesday, June 2, 2027}

\programskip

\pleheading{Plenary}{Main Hall}{Keynote: The Future of Autonomy}{Prof.\ Famous}{Bob Jones}{09:00--10:00}

\programskip

\sesheading{TueA01}{Main Hall}
{Data-Driven Control}
{Alice Smith}{10:30--12:30}

\programskip

\puttalk{TueA01-1\hfill 10.30-10.50}
{\puttalktitle{Learning-Based Model Predictive Control}{paper004}}
{Taylor Speaker}{ETH Zurich}
{}{}
{}{}
{Taylor Speaker}

\programskip

\dayheading{Wednesday, June 3, 2027}

\programskip

\pleheading{Plenary}{Main Hall}{Closing \& Awards Ceremony}{}{Carol Lee}{13:00--14:00}

\programskip
      % your program entries
\programpartindex    % automatic part index
\cpbendprogram

%% --- Abstracts ---
\cpbbeginabstracts
%% Example abstracts file for confprogbook
%%
%% Each abstract is a PDF file included with \includepaper.
%% The first argument is a label (matching the one used in \puttalktitle),
%% and the second is the path to the PDF file.
%%
%% Usage:
%%   \includepaper{paper001}{pdf/paper_001.pdf}
%%   \includepaper{paper002}{pdf/paper_002.pdf}
%%   \includepaper[0.90]{paper003}{pdf/paper_003.pdf}  % custom scale
%%
%% Note: You need actual PDF files for this to compile.
%% Uncomment the lines below and provide your PDF paths.
%%
% \includepaper{paper001}{pdf/paper_001.pdf}
% \includepaper{paper002}{pdf/paper_002.pdf}
% \includepaper{paper003}{pdf/paper_003.pdf}
% \includepaper{paper004}{pdf/paper_004.pdf}
    % \includepaper commands
\cpbendabstracts

%% --- Participants ---
\cpbbeginparticipants
%% Example participant list for confprogbook
%%
%% Usage: \putpart{Prefix}{Name}{Affiliation}{Extra}{Email}
%%
%% Leave unused fields empty: {}

\putpart{}{Jane Researcher}{University of Example}{}{jane@example.edu}
\putpart{}{Bob Collaborator}{Institute of Technology}{}{bob@inst-tech.edu}
\putpart{}{Sam Author}{Technical University}{}{sam@tu.edu}
\putpart{}{Alex Writer}{Research Lab}{}{alex@researchlab.org}
\putpart{}{Taylor Speaker}{ETH Zurich}{}{taylor@ethz.ch}
\putpart{}{Pat Coauthor}{University of Science}{}{pat@unisci.edu}
\putpart{Prof.}{John Doe}{MIT}{Keynote Speaker}{john.doe@mit.edu}
\putpart{Prof.}{Famous Person}{Stanford University}{Keynote Speaker}{}
\putpart{}{Alice Smith}{ETH Zurich}{Organizer}{alice@ethz.ch}
\putpart{}{Bob Jones}{ETH Zurich}{Organizer}{bob.jones@ethz.ch}
\putpart{}{Carol Lee}{ETH Zurich}{Organizer}{carol.lee@ethz.ch}
 % \putpart commands
\cpbendparticipants

%% --- Comments ---
\cpbbegincomments
%% Example comments/organizational info for confprogbook
%%
%% This section is free-form LaTeX. Typical content includes:
%% welcome text, practical information, schedule overview tables,
%% venue maps, sponsor logos, etc.

\section*{Welcome}
The Organizing Committee has the pleasure of welcoming you to the
\emph{10th International Symposium on Examples on Systems and Control},
ETH Zurich, Switzerland.

\section*{Aim}
The aim of this symposium is to promote research activities and to
enhance cooperation between researchers in Systems and Control.

\section*{Scientific Program Overview}
\begin{enumerate}
  \item One plenary lecture by \textit{John Doe} (MIT) on
        \textbf{Advances in Optimal Control}.
  \item One keynote by \textit{Prof.\ Famous} (Stanford) on
        \textbf{The Future of Autonomy}.
  \item Contributed short lectures. See the list of sessions for
        titles and authors.
\end{enumerate}

\section*{Directions for Speakers}
For a contributed lecture, the available time is 20 minutes.
Please leave a few minutes for discussion and room changes.
In each room, a projection screen and HDMI cable are available.

\section*{Registration}
The registration desk is located in the main lobby.

\begin{table}[h]
\centering
\begin{tabular}{|l|c|c|}
\hline
Arrangement & Early & Late \\ \hline
Full registration & \euros{500} & \euros{600} \\ \hline
Student registration & \euros{300} & \euros{400} \\ \hline
Day pass & \euros{200} & \euros{250} \\ \hline
\end{tabular}
\end{table}

\section*{Organization}
The meeting has been organized by
\begin{itemize}
  \item Alice Smith (ETH Zurich)
  \item Bob Jones (ETH Zurich)
  \item Carol Lee (ETH Zurich)
\end{itemize}

\section*{Website}
An electronic version of the Book of Abstracts can be downloaded from the
\href{https://www.example-conference.org}{conference web site}.
     % free-form LaTeX content
\cpbendcomments

\end{document}
\end{lstlisting}

%% ================================================================
\section{Class Options}
\label{sec:options}
%% ================================================================

The \cls{confprogbook} class is built on the standard \cls{book} class.
All options are forwarded directly to \cls{book}. Common choices:

\begin{lstlisting}
\documentclass[a4paper]{confprogbook}      % A4 (default recommended)
\documentclass[letterpaper]{confprogbook}  % US Letter
\documentclass[a4paper,12pt]{confprogbook} % Larger base font
\end{lstlisting}

%% ================================================================
\section{Conference Metadata}
\label{sec:metadata}
%% ================================================================

All metadata commands should be called in the preamble, before
\verb|\begin{document}|. They configure the cover page, headers,
and PDF metadata.

\begin{center}
\begin{tabular}{lp{8cm}}
\toprule
\textbf{Command} & \textbf{Description} \\
\midrule
\cmd{conferencetitle}\marg{text} & Main conference title (shown on cover). \\
\cmd{conferencesubtitle}\marg{text} & Optional subtitle, shown below the title. \\
\cmd{conferenceedition}\marg{text} & Edition string, e.g., \texttt{45th}. \\
\cmd{conferencedate}\marg{text} & Conference dates. \\
\cmd{conferencelocation}\marg{text} & Venue / city / country. \\
\cmd{conferenceshort}\marg{text} & Short name used in page headers. \\
\cmd{booktitle}\marg{text} & Title of the book (default: \textit{Book of Abstracts}). \\
\cmd{editors}\marg{text} & Editor names for the front matter page. \\
\cmd{institution}\marg{text} & Publishing institution name. \\
\cmd{institutionaddress}\marg{text} & Institution address. \\
\cmd{copyrightnotice}\marg{text} & Custom copyright text. \\
\cmd{conferencelogo}\oarg{scale}\marg{file} & Logo image file and optional scale (default 0.2). \\
\cmd{pdfscale}\marg{value} & Scale factor for included abstract PDFs (default 0.95). \\
\bottomrule
\end{tabular}
\end{center}

%% ================================================================
\section{Document Structure Commands}
\label{sec:structure}
%% ================================================================

The class provides high-level commands that set up each major section of
the book. Use them in order inside \verb|\begin{document}...\end{document}|.

\subsection{Cover and Front Matter}

\begin{description}[style=nextline, font=\ttfamily\bfseries, leftmargin=2cm]
  \item[\cmd{makefrontcover}]
    Generates the title page using the metadata you defined. Includes conference
    title, subtitle, dates, location, and book title.

  \item[\cmd{makefrontmatter}]
    Generates the second page with the logo, editor names, book title,
    institution, and copyright notice.
\end{description}

\subsection{Program}

\begin{description}[style=nextline, font=\ttfamily\bfseries, leftmargin=2cm]
  \item[\cmd{cpbbeginprogram}]
    Opens the program section: inserts a ``Part~1'' divider page,
    switches to two-column mode, and activates fancy headers.
    After this, include your program content (day headings, session
    headings, talk entries).

  \item[\cmd{programpartindex}]
    Convenience macro that outputs the standard four-part index
    (Program, Abstracts, Participants, Comments) with dotted leaders.
    Call this after your program content and before \cmd{cpbendprogram}.

  \item[\cmd{cpbendprogram}]
    Closes the program section.
\end{description}

\subsection{Abstracts}

\begin{description}[style=nextline, font=\ttfamily\bfseries, leftmargin=2cm]
  \item[\cmd{cpbbeginabstracts}]
    Opens the abstracts section with a ``Part~2'' divider page and
    switches to one-column mode.

  \item[\cmd{cpbendabstracts}]
    Closes the abstracts section.
\end{description}

\subsection{Participants}

\begin{description}[style=nextline, font=\ttfamily\bfseries, leftmargin=2cm]
  \item[\cmd{cpbbeginparticipants}]
    Opens the participant list section with a ``Part~3'' divider page
    and switches to two-column mode.

  \item[\cmd{cpbendparticipants}]
    Closes the participant list section and starts a new page.
\end{description}

\subsection{Comments / Organizational Info}

\begin{description}[style=nextline, font=\ttfamily\bfseries, leftmargin=2cm]
  \item[\cmd{cpbbegincomments}]
    Opens the comments section (``Part~4'') in one-column mode.
    This section typically contains free-form \LaTeX{} content: welcome text,
    practical info, landscape schedule tables, venue maps, etc.

  \item[\cmd{cpbendcomments}]
    Closes the comments section.
\end{description}

%% ================================================================
\section{Program Entry Commands}
\label{sec:program}
%% ================================================================

These commands typeset the individual entries within the program section.

\subsection{\cmd{dayheading} --- Day Separator}

\begin{lstlisting}
\dayheading{Tuesday, March 24, 2026}
\end{lstlisting}

Produces a full-width colored box (using the ``day'' color scheme) with the
day title centered in bold. Automatically inserts a page break before the heading.

\subsection{\cmd{pleheading} --- Plenary Session Heading}

\begin{lstlisting}
\pleheading{Plenary}{Fagus}{Data-Enabled Predictive Control}
  {Florian Doerfler}{Raphael Jungers}{11:15--12:15}
\end{lstlisting}

Six arguments:
\begin{enumerate}[nosep]
  \item \meta{type} --- e.g., ``Plenary'' (can be empty).
  \item \meta{room} --- room name.
  \item \meta{title} --- talk/event title.
  \item \meta{speaker} --- speaker name (can be empty).
  \item \meta{chair} --- session chair.
  \item \meta{time} --- time slot.
\end{enumerate}

\subsection{\cmd{sesheading} --- Regular Session Heading}

\begin{lstlisting}
\sesheading{TueA01}{Fagus}
  {Neuromorphic Systems and Control I}
  {Rodolphe Sepulchre}{13:15--15:35}
\end{lstlisting}

Five arguments: \meta{code}, \meta{room}, \meta{title}, \meta{chair}, \meta{time}.

\subsection{\cmd{puttalk} --- Individual Talk Entry}

\begin{lstlisting}
\puttalk{TueA01-1\hfill 13.15-13.35}
  {\puttalktitle{Inferring Neuromodulatory Targets}{paper700}}
  {Julien Brandoit}{University of Liege}
  {Pierre Sacre}{University of Liege}
  {Damien Ernst}{University of Liege}
  {Julien Brandoit}
\end{lstlisting}

Nine arguments:
\begin{enumerate}[nosep]
  \item \meta{ID + time} --- talk code and time slot (use \verb|\hfill| to right-align time).
  \item \meta{title line} --- typically a \cmd{puttalktitle} call.
  \item \meta{author1} --- first author name.
  \item \meta{affil1} --- first author affiliation.
  \item \meta{author2} --- second author (empty string \verb|{}| if none).
  \item \meta{affil2} --- second author affiliation.
  \item \meta{author3} --- third author.
  \item \meta{affil3} --- third author affiliation.
  \item \meta{presenter} --- presenting author name.
\end{enumerate}

Leave unused author/affiliation pairs as \verb|{}{}|.

\subsection{\cmd{puttalktitle} --- Talk Title with Page Reference}

\begin{lstlisting}
\puttalktitle{My Talk Title}{paper700}
\end{lstlisting}

Two arguments: \meta{title} and \meta{label}. Produces an italicized title
with dotted leaders to the page where the corresponding abstract PDF begins
(matching the label set by \cmd{includepaper}).

\subsection{\cmd{putplen} --- Part Entry in Program Index}

\begin{lstlisting}
\putplen{Part 1: Programmatic Table of Contents}
  {Overview of scientific program}{P:1}
\end{lstlisting}

Three arguments: \meta{title}, \meta{subtitle}, \meta{label}.

%% ================================================================
\section{Abstract Inclusion}
\label{sec:abstracts}
%% ================================================================

\subsection{\cmd{includepaper} --- Include a Paper PDF}

\begin{lstlisting}
\includepaper{paper700}{pdf/bmsc2026_700.pdf}
\end{lstlisting}

Two mandatory arguments:
\begin{enumerate}[nosep]
  \item \meta{label} --- a unique label (used by \cmd{puttalktitle} for page references).
  \item \meta{filepath} --- path to the PDF file.
\end{enumerate}

Optional argument for custom scale:

\begin{lstlisting}
\includepaper[0.90]{paper700}{pdf/bmsc2026_700.pdf}
\end{lstlisting}

If omitted, the global \cmd{pdfscale} value is used (default: 0.95).
Each included PDF gets the fancy page style and increments the internal
PDF page counter.

%% ================================================================
\section{Participant List}
\label{sec:participants}
%% ================================================================

\subsection{\cmd{putpart} --- Participant Entry}

\begin{lstlisting}
\putpart{}{Alice Smith}{ETH Zurich}{}{alice@ethz.ch}
\end{lstlisting}

Five arguments:
\begin{enumerate}[nosep]
  \item \meta{prefix} --- optional prefix (e.g., ``Prof.'').
  \item \meta{name} --- participant name.
  \item \meta{affiliation} --- institution.
  \item \meta{extra} --- additional info (can be empty).
  \item \meta{email} --- email address (rendered as a \texttt{mailto:} link; leave empty to omit).
\end{enumerate}

%% ================================================================
\section{Color Customization}
\label{sec:colors}
%% ================================================================

The class defines three color groups, each with fill, frame, and text colors.
Customize them in the preamble.

\subsection{Commands}

Each setter takes three arguments. Each argument is a color specification
in the form \verb|{model}{value}|, exactly as you would pass to
\verb|\definecolor|.

\begin{lstlisting}
% Dark blue theme for day headings
\setdaycolors
  {rgb}{0.1, 0.2, 0.5}    % fill
  {rgb}{0.05, 0.1, 0.3}   % frame
  {gray}{1}                % text (white)

% Gold plenary theme
\setplenarycolors
  {rgb}{0.85, 0.65, 0.13}  % fill
  {rgb}{0.72, 0.53, 0.04}  % frame
  {gray}{0}                 % text (black)

% Light gray session theme
\setsessioncolors
  {gray}{0.9}
  {gray}{0.6}
  {gray}{0}

% Table accent column color
\settableaccentcolor{rgb}{0.9, 0.95, 1.0}
\end{lstlisting}

\subsection{Default Colors}

\begin{center}
\begin{tabular}{llll}
\toprule
\textbf{Group} & \textbf{Fill} & \textbf{Frame} & \textbf{Text} \\
\midrule
Day     & \texttt{rgb\{1,.2,.2\}} (red)    & same as fill & \texttt{gray\{0\}} (black) \\
Plenary & \texttt{rgb\{.4,0,.8\}} (purple) & same as fill & \texttt{gray\{1\}} (white) \\
Session & \texttt{rgb\{.8,.8,1\}} (light blue) & same as fill & \texttt{gray\{0\}} (black) \\
\bottomrule
\end{tabular}
\end{center}

%% ================================================================
\section{Landscape Program Tables}
\label{sec:landscape}
%% ================================================================

The comments section often includes landscape overview tables.
The class provides helpers for this:

\begin{description}[style=nextline, font=\ttfamily\bfseries, leftmargin=2cm]
  \item[\cmd{setlandscapecolwidth}\marg{length}]
    Set the width of landscape tables (default: \texttt{244mm}).

  \item[\cmd{rr}]
    Shorthand for \cmd{raggedright}, useful inside \texttt{tabularx} cells.

  \item[\cmd{accentcol}]
    Column specifier prefix for accent-colored columns. Use as:
    \verb|\begin{tabularx}{\cpb@colw}{\accentcol lX|X|...}|
\end{description}

\begin{lstlisting}
\begin{landscape}
\vfill
\centerline{\large\textbf{Tuesday, March 24, 2026}}
\begin{tabularx}{244mm}{\accentcol lX|X|X|Xl|}
\hline
Room & Fagus & Salix & Populus & \\
...
\end{tabularx}
\vfill
\end{landscape}
\end{lstlisting}

%% ================================================================
\section{Utility Commands}
\label{sec:utility}
%% ================================================================

\begin{description}[style=nextline, font=\ttfamily\bfseries, leftmargin=2cm]
  \item[\cmd{euros}\marg{amount}]
    Formats a euro amount: \verb|\euros{600}| renders as ``€~600''.

  \item[\cmd{RED}\marg{text}]
    Highlights text in the highlight color (default red).

  \item[\cmd{programskip}]
    Shorthand for \verb|\medskip\nopagebreak|, the most common spacing
    pattern between program entries.
\end{description}

%% ================================================================
\section{Customizing Part Names}
\label{sec:partnames}
%% ================================================================

If you want to rename the four standard parts (e.g., to a different language),
use:

\begin{lstlisting}
\setpartname{I}{Partie 1}{Table des matieres programmatique}
\setpartname{II}{Partie 2}{Resumes}
\setpartname{III}{Partie 3}{Liste des participants}
\setpartname{IV}{Partie 4}{Informations pratiques}
\end{lstlisting}

Note: this currently sets internal macros. The high-level structure
commands (\cmd{cpbbeginprogram}, etc.) use hardcoded English strings for
the divider pages. For full localization, override the structure commands
or redefine them after loading the class.

%% ================================================================
\section{Package Dependencies}
\label{sec:dependencies}
%% ================================================================

The class automatically loads the following packages. You do \emph{not}
need to load them yourself.

\begin{center}
\begin{tabular}{ll}
\toprule
\textbf{Package} & \textbf{Purpose} \\
\midrule
\pkg{amsmath}, \pkg{amsfonts}, \pkg{amssymb}, \pkg{amstext} & Math support \\
\pkg{graphicx} & Image inclusion \\
\pkg{xcolor} (with \texttt{table} option) & Color support \\
\pkg{latexsym} & Additional symbols \\
\pkg{ifthen}, \pkg{etoolbox} & Conditionals \\
\pkg{tabularx}, \pkg{ltablex} & Flexible tables \\
\pkg{lscape} & Landscape pages \\
\pkg{fancyhdr} & Custom headers/footers \\
\pkg{pdfpages} & PDF inclusion \\
\pkg{fontenc} (T1) & Font encoding \\
\pkg{hyperref} & Hyperlinks and PDF metadata \\
\pkg{pgfkeys} & Key-value processing \\
\pkg{xparse} & Extended command definitions \\
\bottomrule
\end{tabular}
\end{center}

%% ================================================================
\section{Migrating from a Manual Setup}
\label{sec:migration}
%% ================================================================

If you have an existing conference book that uses manual \cmd{newcommand}
definitions (like the original Benelux Meeting setup), migration is
straightforward:

\begin{enumerate}
  \item Replace \verb|\documentclass{book}| with
        \verb|\documentclass[a4paper]{confprogbook}|.
  \item Delete your package loading block --- the class loads everything.
  \item Delete your custom command definitions (\cmd{dayheading},
        \cmd{pleheading}, \cmd{sesheading}, etc.) --- they are now
        provided by the class.
  \item Delete your color definitions --- the class provides defaults.
        Customize with \cmd{setdaycolors} etc.\ if needed.
  \item Delete your page layout settings --- the class handles this.
  \item Add metadata commands in the preamble (\cmd{conferencetitle}, etc.).
  \item Replace your manual section structure with the high-level
        commands (\cmd{makefrontcover}, \cmd{cpbbeginprogram}, etc.).
  \item Replace long \verb|\includepdf[...]| calls with
        \verb|\includepaper{label}{file}|.
  \item Your \cmd{puttalk}, \cmd{putpart}, and content files
        work \emph{unchanged}.
\end{enumerate}

%% ================================================================
\section{Frequently Asked Questions}
\label{sec:faq}
%% ================================================================

\paragraph{Can I use this with \LuaLaTeX{} or \XeLaTeX{}?}
Yes. The class uses standard packages that work across all major engines.
If using \LuaLaTeX/\XeLaTeX, you may want to add \pkg{fontspec} in your
document preamble for system fonts.

\paragraph{Can I add extra parts beyond the four standard ones?}
Yes. The structure commands are just convenience macros. You can freely
add \verb|\onecolumn|, \verb|\twocolumn|, \verb|\cleardoublepage|, etc.\
and create additional sections as needed.

\paragraph{Can I use this on Overleaf?}
Yes. Upload \texttt{confprogbook.cls} to your Overleaf project root
(same directory as your \texttt{main.tex}) and it will be found
automatically. See Section~\ref{sec:overleaf} for details.

\paragraph{How do I change the header text?}
Use \cmd{conferenceshort} for the left/right header and \cmd{booktitle}
for the other side.

%% ================================================================
\section{Publishing to CTAN and Overleaf}
\label{sec:publishing}
%% ================================================================

\subsection{CTAN Submission}
\label{sec:ctan}

To make the class available to the entire \LaTeX{} community via
CTAN (Comprehensive \TeX{} Archive Network):

\begin{enumerate}
  \item \textbf{Prepare the TDS-compliant archive:}
  \begin{verbatim}
  confprogbook/
    confprogbook.cls
    confprogbook-doc.tex
    confprogbook-doc.pdf
    README.md
    LICENSE
    example/
      main.tex
      program.tex
      abstracts.tex
      participants.tex
      comments.tex
  \end{verbatim}

  \item \textbf{Compile the documentation:}
  \begin{verbatim}
  pdflatex confprogbook-doc.tex
  pdflatex confprogbook-doc.tex
  \end{verbatim}

  \item \textbf{Create a zip archive:}
  \begin{verbatim}
  zip -r confprogbook.zip confprogbook/
  \end{verbatim}

  \item \textbf{Submit to CTAN:}
  \begin{itemize}[nosep]
    \item Go to \url{https://ctan.org/upload}.
    \item Fill in the metadata: name (\texttt{confprogbook}), version,
          summary, license (LPPL 1.3c), author.
    \item Upload the zip file.
    \item CTAN volunteers will review and integrate the package
          (typically within a few days).
  \end{itemize}

  \item \textbf{TeX Live / MiKTeX integration:}
    Once on CTAN, the package will be picked up by TeX~Live and MiKTeX
    in their next update cycle. Users can then install it via their
    package manager.
\end{enumerate}

\subsection{Overleaf Availability}
\label{sec:overleaf}

Overleaf uses TeX~Live, so once your package is on CTAN and included in
TeX~Live, it becomes automatically available on Overleaf. There are also
faster paths:

\begin{enumerate}
  \item \textbf{Immediate use (any project):} Upload \texttt{confprogbook.cls}
        directly into your Overleaf project.

  \item \textbf{Overleaf template gallery:} Create a template project on
        Overleaf with the class and example files, then submit it to the
        \href{https://www.overleaf.com/latex/templates}{Overleaf Gallery}
        via \texttt{Submit to Gallery} in the project menu.

  \item \textbf{Automatic via CTAN $\to$ TeX~Live:} After CTAN acceptance
        and the next TeX~Live update, users can use
        \verb|\documentclass{confprogbook}| directly without uploading
        the \texttt{.cls} file.
\end{enumerate}

%% ================================================================
\section{Version History}
\label{sec:history}
%% ================================================================

\begin{description}
  \item[v1.0 (2026/03/11)] Initial release. All core functionality:
    metadata commands, structure commands, program entries, abstract
    inclusion, participant lists, color customization.
\end{description}

%% ================================================================
\section{Acknowledgments}
\label{sec:acknowledgments}
%% ================================================================

This class was inspired by the \LaTeX{} setup used for the
\emph{Benelux Meeting on Systems and Control} Book of Abstracts,
originally developed by the conference organizers. Special thanks to
Gianluca Bianchin, Julien Hendrickx, Rapha\"el Jungers, and Erjen Lefeber.

\end{document}
